\chapter{Gentamicin Level}

\section*{What Is This Test?}
The gentamicin level measures the serum concentration of gentamicin, an aminoglycoside antibiotic used to treat serious Gram-negative bacterial infections (particularly Pseudomonas aeruginosa, Enterobacteriaceae including Escherichia coli, Klebsiella, Enterobacter, Serratia, and Proteus species) and as synergistic therapy for serious Gram-positive infections (enterococcal and staphylococcal endocarditis, in combination with a beta-lactam antibiotic). Gentamicin binds to the 30S ribosomal subunit, inhibiting bacterial protein synthesis and causing misreading of the genetic code, resulting in a bactericidal effect. Gentamicin is administered intravenously or intramuscularly and has a narrow therapeutic index. Toxicity is concentration-dependent and cumulative: nephrotoxicity (acute tubular necrosis, ATN) occurs in 10--25\% of patients and is typically reversible with drug discontinuation; ototoxicity (irreversible vestibular and cochlear damage --- high-frequency hearing loss progressing to lower frequencies, and bilateral vestibular dysfunction with oscillopsia, ataxia) is irreversible and occurs in 5--15\% of patients. The incidence of both nephrotoxicity and ototoxicity is directly related to sustained elevated trough concentrations. Therapeutic drug monitoring (TDM) measures peak (30 minutes after a 30-minute IV infusion, or 60 minutes after IM injection) and trough (immediately before the next dose) levels. Extended-interval (once-daily) dosing (5--7 mg/kg every 24 hours) has largely replaced traditional multiple-daily dosing (1.5--2 mg/kg every 8 hours) because it achieves higher peak/MIC ratios (improving bacterial killing via concentration-dependent killing) and lower trough levels (reducing nephrotoxicity and ototoxicity) \cite{begg1995suggested}.

\section*{Alternative Names}
Gentamicin Level; Gentamicin Peak and Trough; Aminoglycoside Level; Gentamicin TDM; Serum Gentamicin; Garamycin Level

\section*{Principle}
Gentamicin is measured by particle-enhanced turbidimetric inhibition immunoassay (PETINIA) or enzyme-multiplied immunoassay technique (EMIT) on automated clinical chemistry analyzers. Patient serum containing gentamicin is incubated with anti-gentamicin antibody and gentamicin-coated microparticles. The rate of particle aggregation inhibition is proportional to gentamicin concentration.

\section*{History}
Gentamicin was discovered in 1963 from the soil actinomycete Micromonospora purpurea by Weinstein and colleagues at Schering Corporation. It became available in the United States in 1969. The introduction of once-daily dosing regimens in the 1990s was based on the recognition of two key pharmacodynamic concepts: concentration-dependent killing (higher peak/MIC ratio increases bacterial killing) and the post-antibiotic effect (suppression of bacterial growth that persists after antibiotic concentrations fall below the MIC) \cite{freeman1997once}. The Hartford nomogram for extended-interval dosing was developed in 1999 \cite{nicolau1995experience}.

\section*{Why This Test Is Ordered}
Therapeutic drug monitoring to maximize efficacy and minimize nephrotoxicity and ototoxicity. Gentamicin levels should be monitored in: patients with renal impairment or rapidly changing renal function, prolonged therapy (>3--5 days), serious infections requiring TDM-optimized dosing, critically ill patients, patients with large volume of distribution changes (ascites, burns, pregnancy, cystic fibrosis), and patients receiving concurrent nephrotoxic drugs.

\section*{Primary vs Secondary Test}
This is a primary therapeutic drug monitoring test.

\section*{How This Test Is Performed}
Blood sample in a serum separator tube (SST, gold-top) or red-top tube. Peak level: 30 minutes after the end of a 30-minute IV infusion (or 60 minutes after IM injection). Trough level: immediately before the next dose. For extended-interval (once-daily) dosing: serum concentration is measured 6--14 hours after the dose, and the result is plotted on the Hartford nomogram to determine the dosing interval (every 24 hours, every 36 hours, or every 48 hours).

\section*{What Preparation Is Needed}
No fasting is required. Accurate recording of the dose and infusion times is essential because interpretation depends on the exact time of collection relative to dosing. The peak level must be drawn 30 minutes after the end of a 30-minute intravenous infusion (or 60 minutes after intramuscular injection) and the trough immediately before the next dose; for extended-interval dosing, a single level is drawn 6--14 hours after the dose. Blood should be drawn from a site other than the line used for gentamicin infusion to avoid sample contamination. The laboratory needs to know the dosing regimen (traditional or extended-interval) and the time of the last dose \cite{avent2011current}.

\section*{Contraindications and Precautions}
\begin{itemize}
  \item There are no absolute contraindications to venipuncture for gentamicin level measurement. Interpretive cautions:
  \item (1) the Hartford nomogram for extended-interval dosing is not validated in patients with endocarditis, pregnancy, burns, cystic fibrosis, ascites, or creatinine clearance below 20 mL/min, and should not be used in these settings.
  \item (2) Concurrent nephrotoxic drugs (loop diuretics, NSAIDs, amphotericin B, vancomycin, cisplatin) increase the risk of nephrotoxicity.
  \item (3) A single trough level is not interpretable without knowing whether dosing is traditional or extended-interval, because the target troughs differ.
  \item (4) Levels obtained from the infusion line or at the wrong time are unreliable and should be redrawn.
\end{itemize}
\section*{Risks}
Minimal risk. Slight pain or bruising at the venipuncture site. Rarely: hematoma, infection, or vasovagal syncope.

\section*{What Happens After the Test}
Levels are typically reported within a few hours. The result is interpreted against the target peak and trough for the dosing regimen (traditional or extended-interval) and the patient's renal function. A trough above target prompts extension of the dosing interval or dose reduction, whereas a low peak prompts a dose increase; after any adjustment, repeat levels are obtained at the next steady state \cite{bartal2003pharmacokinetic}. Renal function and electrolytes are monitored during continued therapy, and audiometry is performed in high-risk patients.

\section*{Understanding Results}

\subsection*{Below Therapeutic Range (Subtherapeutic)}
Peak <4--5 mcg/mL (traditional multiple-daily dosing). Inadequate dosing or augmented renal clearance. Risk of suboptimal bacterial killing, treatment failure, and emergence of resistance (adaptive resistance develops with prolonged low-level gentamicin exposure).

\subsection*{Above Therapeutic Range (Potentially Toxotic)}
Trough >2 mcg/mL (traditional multiple-daily dosing) or >0.5--1 mcg/mL (extended-interval dosing). Increased risk of nephrotoxicity (non-oliguric ATN presenting with rising serum creatinine after 5--7 days of therapy, usually reversible with drug discontinuation; hypomagnesemia and hypokalemia from renal tubular toxicity may precede the rise in creatinine). Ototoxicity (cumulative dose-dependent; vestibular toxicity presents with oscillopsia, ataxia, and disequilibrium; high-frequency hearing loss is often subclinical and detected by audiometry before the patient notices symptoms).

\section*{Normal Results}
Traditional multiple-daily dosing: Peak 4--10 mcg/mL (severe Gram-negative infections); 3--5 mcg/mL (synergy for enterococcal endocarditis). Trough <2 mcg/mL (<1 mcg/mL preferred). Extended-interval (once-daily) dosing: Peak 15--25 mcg/mL (30 minutes after a 30-minute IV infusion). Trough <0.5--1 mcg/mL. The Hartford nomogram: measure a single serum level 6--14 hours after the dose; plot to determine the dosing interval. The Hartford nomogram should NOT be used for patients with endocarditis, pregnancy, burns, cystic fibrosis, ascites, or CrCl <20 mL/min \cite{winter2010basic}.

\section*{Follow-up Tests}
Serum creatinine and BUN at baseline and at least 2--3 times weekly during therapy. Serum electrolytes (magnesium, potassium) --- gentamicin causes renal magnesium and potassium wasting. Audiometry (high-frequency pure-tone audiometry) for high-risk patients or prolonged therapy. Repeat peak and trough levels with dose adjustment as indicated.

\section*{References}
\bibliographystyle{unsrtnat}
\bibliography{references}

\clearpage
\section*{Regulatory Status and Authorization Date}
This chapter describes a test category, not a specific commercial product. FDA, CDSCO, EU IVDR, and MHRA approval or registration dates therefore cannot be assigned without the assay manufacturer, product name, instrument, and intended use. For laboratory-developed tests, an individual agency approval date may not exist. See the Preface for the interpretation of regulatory status.

\clearpage
